\documentclass[10pt]{article}

% Industry resume, built from the academic CV template.
% margins: 0.5in all around
\usepackage[margin=0.5in]{geometry}
\pagestyle{empty}

\usepackage{marvosym}
\usepackage[plainpages=false, colorlinks=true,
   citecolor=blue, filecolor=blue, linkcolor=blue,
   urlcolor=blue]{hyperref}
\usepackage{enumitem}
\setlist[itemize]{leftmargin=1.3em, itemsep=0pt, topsep=1pt, parsep=0pt}

% Section titles with a thin rule underneath (same style as the CV)
\usepackage{titlesec}
\titleformat{\section}{\large\bfseries}{}{0pt}{}[\vspace{1pt}\titlerule]
\titlespacing*{\section}{0pt}{6pt}{3pt}

% One experience entry: title, organization, location, dates
\newcommand{\entry}[4]{%
  \noindent{\bf #1}, #2 {\it (#3)} \hfill #4\par\vspace{1pt}}

\begin{document}
\emergencystretch=2em

%#################################################################################
% HEADER
\begin{center}
  {\bfseries \LARGE Irabiel Romero} \\[0.06in]
  Applied Mathematics Ph.D. (Dec.\ 2026) \textbar\ Scientific Computing \& Machine Learning \\[0.04in]
  \Mobilefone\ (661) 240-7760 \quad
  \Letter\ \href{mailto:irabielr96@gmail.com}{irabielr96@gmail.com} \quad
  \href{https://www.linkedin.com/in/irabiel}{linkedin.com/in/irabiel} \quad
  \href{https://irabiel.github.io}{irabiel.github.io} \quad
  Merced, CA
\end{center}

%#################################################################################
\section*{Summary}
\noindent
Applied mathematician building scalable numerical methods for estimation and
optimization under uncertainty: Bayesian inverse problems, optimal sensor
placement, large-scale nonlinear and PDE-constrained optimization, multigrid
solvers, and ML surrogates inside constrained optimizers. Three research
programs at Lawrence Livermore National Laboratory.

%#################################################################################
\section*{Education}
\noindent
{\bf Ph.D., Applied Mathematics}, University of California, Merced \hfill Expected Dec.\ 2026\\
{\it Thesis:} Scalable Methods for Estimation of Power Systems under Uncertainty
(Advisors: Noemi Petra, Roummel Marcia)\\[2pt]
{\bf M.S., Applied Mathematics}, California State University, Long Beach \hfill May 2021\\
{\it Thesis:} Numerical Methods for Multiscale Modeling in Bacterial Colony Growth
(Outstanding Thesis Award)\\[2pt]
{\bf B.S., Applied Mathematics} and {\bf B.A., Physics}, California State
University, Long Beach \hfill May 2018

%#################################################################################
\section*{Experience}

\entry{Graduate Research Intern}{CASC, Lawrence Livermore National Laboratory}%
{Livermore, CA}{Jun.\ 2026 -- Aug.\ 2026}
\begin{itemize}
\item Integrated a machine learning surrogate into a large-scale nonlinear
  optimization solver, replacing stability constraints that are
  computationally intractable to enforce directly.
\item Designed and ran validation tests comparing surrogate-constrained
  solutions against full simulation.
\end{itemize}

\entry{Graduate Student Researcher}{Lawrence Livermore National Laboratory}%
{Remote}{Jan.\ 2025 -- May 2026}
\begin{itemize}
\item Extended LLNL's security-constrained AC optimal power flow solver
  (ExaJuGO, Julia) to analyze wildfire-driven multi-failure contingencies on
  the California Test System; built interactive visualization tools.
\item Developed a CNN surrogate for generator stability and embedded it in
  the optimizer to enforce transient stability constraints that would
  otherwise be computationally intractable.
\end{itemize}

\entry{Graduate Student Researcher}{University of California, Merced}%
{Merced, CA}{Jan.\ 2023 -- Present}
\begin{itemize}
\item Formulated a Bayesian inverse problem to estimate system states and
  uncertain model parameters from sparse sensor data, and an optimal
  experimental design method for sensor placement.
\item Developed a generic algebraic multigrid solver for M-matrices and a
  PDE-constrained inverse problem (advection--diffusion) that recovers a
  diffusion parameter from sparse data.
\item Led an LLNL project with CASC and the UQ group (Fall 2026): built an ML
  surrogate that learns the distribution of stability outcomes over all
  fault locations to estimate the probability of instability, with conformal
  prediction giving calibrated, distribution-free bounds on that probability.
\end{itemize}

\entry{Data Science Intern}{Lawrence Livermore National Laboratory}%
{Livermore, CA}{Summers 2023, 2024}
\begin{itemize}
\item Data Science Summer Institute (2024): built a parallel framework (C++,
  MPI, HDF5) using particle swarm optimization to generate large simulation
  databases; trained CNNs to classify heartbeat irregularities from ECG
  signals.
\item Data Science Challenge (2023): led a team developing a 1D CNN that
  reconstructs the heart's transmembrane voltages and activation times from
  ECG signals.
\end{itemize}

\entry{Graduate Researcher}{CSU Long Beach \& UC San Diego}%
{Long Beach and San Diego, CA}{Jun.\ 2018 -- May 2021}
\begin{itemize}
\item Developed a Newton--multigrid solver for stiff nonlinear
  reaction--diffusion PDEs, coupled with a kinetic Monte Carlo simulation.
\end{itemize}

%#################################################################################
\section*{Selected Publications}
\begin{itemize}
\item {\bf I. Romero} et al. ``Wildfire-driven N-k contingency analysis for
  security-constrained AC optimal power flow in the California power grid.''
  IEEE EPEC, 2026.
\item {\bf I. Romero}, R. Marcia, I. Aravena, N. Petra. ``Optimal PMU
  placement for state estimation with grid parameter uncertainty.'' IEEE PES
  Grid Edge, 2025.
\item A. Ranganath, {\bf I. Romero}, M. Singhal, R. Marcia. ``Quasi-Adam:
  Accelerating Adam using quasi-Newton approximations.'' IEEE ICMLA, 2024.
\end{itemize}

%#################################################################################
\section*{Skills \& Awards}
\noindent
\begin{tabular}{@{}l@{\hspace{0.8em}}p{0.78\textwidth}}
Languages \& HPC: & Python, C++, Julia, MATLAB; MPI, OpenMP, HDF5, Git \\
Methods: & PDE-constrained optimization, inverse problems, Bayesian inference, UQ, conformal prediction \\
ML \& libraries: & PyTorch, TensorFlow, JAX, SciPy, IPOPT, FEniCS, hIPPYlib; CNNs, surrogate modeling \\
Awards: & SIAM NCC poster awards (2025, 2026); Fields Family Fellowship; Outstanding Thesis Award
\end{tabular}

%#################################################################################
\end{document}
